
This study began with a coverage paradox: pylint flags 100\% of HumanEval baseline failures, yet pylint-guided repair reduces HumanEval pass@1 by 4.3 percentage points. The paradox resolves when coverage is decomposed by flag category. The 100\% total coverage is driven by Convention-category style rules (C0304: missing final newline, C0114: missing module docstring) that fire universally on LLM-generated code snippets regardless of functional correctness. Only 12.5\% of HumanEval failures receive a functional Error or Warning category flag; mypy coverage is 0\%. When the LLM receives style feedback as its primary repair signal for an algorithmic failure, it may produce a stylistically cleaner but logically incorrect solution.

The three main contributions are: (1) the first iso-compute comparison confirming that execution feedback significantly dominates pylint/mypy on HumanEval and MBPP (McNemar p=0.0001 and p<$10^{-18}$ respectively); (2) the style-function dissociation measurement (94.3\% C-category, 12.5\% E+W functional coverage) as a novel decomposition of why pylint coverage overstates feedback informativeness; and (3) the benchmark asymmetry finding that static analysis repair utility is moderated by task complexity.

\textbf{Future directions.} Qwen2.5-Coder-7B replication to assess model generalizability; per-problem qualitative analysis for regression cases under pylint repair; W+E-only pylint filtering to evaluate whether removing C-category style noise improves pylint repair efficacy; larger budget comparison (B=2000+) for multi-round characterization; extension to code-specialized models and non-Python benchmarks.

For LLM code repair, the relevant question is not whether a static analysis tool detects \textit{some} issue in a failing solution --- it is whether the detected issue corresponds to the functional reason for the failure. Pylint detects style conventions; execution detects functional failures. At B=1000 on functional correctness benchmarks, the difference in repair efficacy is +40.2pp on MBPP.

